\section{The actual defect from signed facet cofactors}\label{sec:cofactor}
Let $K$ be a full-dimensional $d$-polytope with exactly $N=d+2$
irredundant facets. Put $n=d+1$. For facet $i$, let $\nu_i$ be its outward
unit normal, $s_i>0$ its actual area, and $h_i$ its support number. Set
\[
 g_i=s_i\nu_i,\qquad U_i=(g_i,s_ih_i)\in\R^{d+1},\qquad V=|K|>0.
\]
The divergence identities give
\begin{equation}\label{eq:cofactor-sum}
 \sum_i g_i=0,\qquad \sum_i s_ih_i=dV,
 \qquad\sum_iU_i=(0,\ldots,0,dV).
\end{equation}
The horizontal rows of $U$ have rank $d$ because the facet normals span
$\R^d$. The last column-sum identity supplies an independent vertical
direction, so $U$ has rank $n$. Define the signed maximal cofactors
\begin{equation}\label{eq:cofactor-z}
 z_i=(-1)^{i-1}\det(U_1,\ldots,U_{i-1},U_{i+1},\ldots,U_N).
\end{equation}
Then $z\ne0$ and $\sum_i z_iU_i=0$. Translating the origin to the interior
of $K$ leaves these determinants invariant by a row operation and makes
all $h_i>0$. Thus $\sum_i z_is_ih_i=0$ forces both positive and negative
entries in $z$.

\begin{theorem}[Exact cofactor defect]\label{thm:cofactor}
With $L=\sum_i|z_i|$ and $T=\sum_{i<j}|z_i-z_j|$, the original invariant
satisfies
\begin{equation}\label{eq:cofactor-invariant}
 a(K)=\frac LT,\qquad
 \delta(K):=n-\frac1{a(K)}=n-\frac TL,\qquad
 e(K)=\frac{\delta(K)}{n(n-\delta(K))}.
\end{equation}
In particular,
\begin{equation}\label{eq:cofactor-min}
 \delta(K)=\frac2L\sum_{\substack{i<j\\z_iz_j>0}}
                            \min(|z_i|,|z_j|).
\end{equation}
\end{theorem}
\begin{proof}
Let $H=\sum_{i<j}|\det(g_k:k\notin\{i,j\})|$, with the remaining
columns in increasing order. The actual finite-facet identity proved
in Appendix~\ref{prop:actual-moments} gives $a(K)=L/(dVH)$. We show
\begin{equation}\label{eq:cofactor-difference}
 dV\,|\det(g_k:k\notin\{i,j\})|=|z_i-z_j|.
\end{equation}
Fix $i<j$. In the $n$-column matrix omitting column $i$, replace its
column $U_j$ by $\sum_kU_k$. By multilinearity only the original term
and the term replacing $U_j$ by $U_i$ survive. Reordering the latter
columns with the signs in \eqref{eq:cofactor-z} gives $z_i-z_j$ up to an
overall sign. The replacement column is $(0,\ldots,0,dV)$ by
\eqref{eq:cofactor-sum}, and expansion along it gives the other side of
\eqref{eq:cofactor-difference}, again up to sign. Hence $dVH=T$ and
$a=L/T$. Solving $e=a-1/n$ proves \eqref{eq:cofactor-invariant}; all
its denominators are positive.

For real $x,y$, the difference $|x|+|y|-|x-y|$ equals
$2\min(|x|,|y|)$ when $xy>0$, and equals zero otherwise. Each $|z_i|$
occurs in $n$ unordered pairs. Summing this identity gives
\eqref{eq:cofactor-min}.
\end{proof}

\subsection{Signed concentration}
Sort the positive entries as $p_1\ge\cdots\ge p_p>0$ and the absolute
values of the negative entries as $q_1\ge\cdots\ge q_q>0$. Then
\begin{equation}\label{eq:cofactor-sorted}
 \delta(K)=\frac2L\left(\sum_{i=2}^p(i-1)p_i+
                              \sum_{j=2}^q(j-1)q_j\right).
\end{equation}
For $\varepsilon=(L-p_1-q_1)/L$ this gives
\begin{equation}\label{eq:cofactor-concentration}
 2\varepsilon\le\delta(K)\le2\max(p-1,q-1)\varepsilon
                          \le2d\varepsilon.
\end{equation}
Indeed each omitted weight in \eqref{eq:cofactor-sorted} has coefficient
between one and $\max(p-1,q-1)$, and both sign classes are nonempty.
At least $1-\delta(K)/2$ of the normalized absolute cofactor mass is
therefore concentrated on one positive and one negative entry. The
weighted identity, rather than only this comparison, determines the defect.
Common nonzero rescaling and overall sign reversal leave all formulas
unchanged. Zero cofactors, including pyramid degeneracies, are permitted;
no simplicity or genericity condition is required. Arbitrary signed
vectors are not asserted to be realizable facet data.

\subsection{Explicit clipping coordinates}
Suppose $K=P\cap\{f\ge0\}$ with $f=\sum_if_i\lambda_i$, and now assume
that all $n$ old facets and the cut facet remain irredundant. Use the
partition $f_i=-a_i<0$ on $I$, $f_j=b_j>0$ on $J$, allowing zero values.
For any origin there are unique positive scales $r_i,q$ such that
\[
 U_i=r_i(-\nabla\lambda_i,\lambda_i(0)),\qquad
 U_c=q(-\nabla f,f(0)).
\]
Since $\sum_i\lambda_i=1$, the column-sum identity and independence of the
affine barycentric functions imply
\begin{equation}\label{eq:cofactor-scales}
 r_i+qf_i=R,\qquad R=dV>0,
\end{equation}
so each $r_i=R-qf_i$ is positive. The affine identity for $f$ gives the
kernel vector, up to common scale,
\[
 z_c=1,\qquad z_i=-qf_i/r_i.
\]
Write
\[
 \alpha_i=\frac{qa_i}{R+qa_i}\in(0,1),\qquad
 \beta_j=\frac{qb_j}{R-qb_j}>0.
\]
The cut entry and each $\alpha_i$ have the same sign, with
$\min(1,\alpha_i)=\alpha_i$, so \eqref{eq:cofactor-min} becomes
\begin{equation}\label{eq:cofactor-clipping}
 \delta(K)=
 \frac{2\left[\sum_i\alpha_i+\sum_{i<i'}\min(\alpha_i,\alpha_{i'})+
                         \sum_{j<j'}\min(\beta_j,\beta_{j'})\right]}
 {1+\sum_i\alpha_i+\sum_j\beta_j}.
\end{equation}
Zero $f_i$ give zero cofactors. If clipping removes an entire old facet,
that redundant facet must be deleted before applying this calculation;
one may not divide by a vanished area.

For the equal-intercept truncation
\begin{equation}\label{eq:truncation-body}
 K_t=\{x_i\ge0:\ t\le\textstyle\sum_i x_i\le1\},\qquad0<t<1,
\end{equation}
the normalized cofactor vector is
\[
 z_c=1,\qquad z_0=t^d,\qquad z_1=\cdots=z_d=-\beta,
 \qquad\beta=\frac{t^{d-1}(1-t)}{1-t^{d-1}}.
\]
The old opposite facet retains its whole area, each other old facet
retains the fraction $1-t^{d-1}$, and the cut area scale is the fraction
$t^{d-1}$. These facts, or substitution in \eqref{eq:cofactor-scales},
give the displayed vector. Therefore
\begin{equation}\label{eq:truncation-delta}
 \delta(K_t)=\frac{2[t^d+\binom d2\beta]}{1+t^d+d\beta}.
\end{equation}
Substitution into \eqref{eq:cofactor-invariant} yields the exact original
defect and its fixed-dimensional leading term
\begin{equation}\label{eq:truncation-exact}
 \begin{split}
 e(K_t)&=\frac{t^{d-1}[d(d-1)-(d+1)(d-2)t-2t^d]}
 {(d+1)(1-t^d)(d+1+(d-1)t^{d-1})},\\
 e(K_t)&=\frac{d(d-1)}{(d+1)^2}t^{d-1}+O(t^d).
 \end{split}
\end{equation}
For a square followed by $d-2$ pyramids, the four nonzero cofactors are
proportional to $(1,-1,1,-1)$ and the remaining $d-2$ vanish. Thus
$\delta=1$ and $e=1/[d(d+1)]$. The direct facet certificate in the next
section checks the same actual invariant without assuming realizability
of a merely formal cofactor vector.
